% Luc Destombe --- licence CC BY-NC-SA 4.0
% https://creativecommons.org/licenses/by-nc-sa/4.0/deed.fr
% Fichier autonome : compiler avec pdflatex, deux fois (signets, nombre de pages).
\def\lyceeversion{1.0}
\documentclass[a4paper,10pt]{article}
\makeatletter
\newif\iflycee@niveaux
\lycee@niveauxtrue
\newif\iflycee@palatino
\lycee@palatinotrue

\RequirePackage[utf8]{inputenc}
\RequirePackage[T1]{fontenc}
\RequirePackage[french]{babel}
\RequirePackage{etoolbox}

\newcommand{\ShorthandsOff}{\@ifundefined{shorthandoff}{}{\shorthandoff{;:!?}}}
\newcommand{\ShorthandsOn}{\@ifundefined{shorthandon}{}{\shorthandon{;:!?}}}
\ShorthandsOff
\AtBeginDocument{\ShorthandsOn}
\AtBeginEnvironment{tikzpicture}{\ShorthandsOff}

\RequirePackage{geometry}
\RequirePackage{fancyhdr}
\RequirePackage{lastpage}
\RequirePackage{multicol}
\RequirePackage{array}
\RequirePackage{enumitem}
\RequirePackage{xcolor}
\RequirePackage{needspace}

\RequirePackage{amsmath,amssymb}
\RequirePackage{mathpazo}
\renewcommand{\ttdefault}{lmtt}
\DeclareMathSizes{10.95}{10.95}{8}{5.5}
\begingroup\catcode`\_=\active \catcode`\^=\active
\gdef_#1{\sb{\scriptscriptstyle #1}}
\gdef^#1{\sp{\scriptscriptstyle\lycee@moinsepais #1}}
\endgroup
\newcommand\lycee@moins{\mathbin{%
  \setbox\z@\hbox{$\m@th\scriptscriptstyle\lycee@moinsorig$}%
  \dimen@=.55\wd\z@ \advance\dimen@-.5pt
  \hbox to.75\wd\z@{\hss$\m@th\scriptscriptstyle\vcenter{\hbox{%
    \tikz\draw[line width=.5pt,line cap=round](0,0)--(\the\dimen@,0);}}$\hss}}}
\begingroup\lccode`\~=`\- \lowercase{\endgroup\let~}\lycee@moins
\newcommand\lycee@moinsepais{\mathcode`\-="8000 }
\AtBeginDocument{\mathchardef\lycee@moinsorig=\mathcode`\-\relax}
\AtBeginDocument{\catcode`\_=12 \mathcode`\_="8000
  \catcode`\^=12 \mathcode`\^="8000 }
\newcommand\lycee@exposants{%
  \fontdimen13\textfont2=.48\fontdimen6\textfont2
  \fontdimen14\textfont2=.48\fontdimen6\textfont2
  \fontdimen15\textfont2=.48\fontdimen6\textfont2
  \fontdimen13\scriptfont2=.48\fontdimen6\scriptfont2
  \fontdimen14\scriptfont2=.48\fontdimen6\scriptfont2
  \fontdimen15\scriptfont2=.48\fontdimen6\scriptfont2}
\every@math@size\expandafter{\the\every@math@size\lycee@exposants}
\newlength{\retraitcalculs}
\setlength{\retraitcalculs}{2em}

\RequirePackage{tikz}
\usetikzlibrary{arrows.meta,calc,babel,shapes.misc}
\RequirePackage[most]{tcolorbox}
\tcbuselibrary{skins,breakable}

\definecolor{coulDef}{HTML}{1F4E79}
\definecolor{coulProp}{HTML}{7030A0}
\definecolor{coulRem}{HTML}{C55A11}
\definecolor{coulEx}{HTML}{2E7D32}
\definecolor{coulExo}{HTML}{B03A2E}
\definecolor{coulPart}{HTML}{1F4E79}
\definecolor{coulMeth}{HTML}{1B6E4F}
\definecolor{coulCourbe}{HTML}{0B5ED7}
\definecolor{coulCourbeB}{HTML}{C00000}
\definecolor{coulCourbeC}{HTML}{2E7D32}
\definecolor{coulGrille}{HTML}{8C8C8C}
\definecolor{coulRep}{HTML}{1B6E4F}
\definecolor{coulBar}{HTML}{2E7D32}

\newcommand{\R}{\mathbb{R}}
\newcommand{\Rs}{\mathbb{R}^{*}}
\renewcommand{\le}{\leqslant}
\renewcommand{\ge}{\geqslant}
\renewcommand{\approx}{\simeq}
\newcommand{\vect}[1]{\overrightarrow{#1}}
\newcommand{\norme}[1]{\left\lVert #1 \right\rVert}
\newcommand{\intervalle}[2]{\left[#1\,;#2\right]}
\RequirePackage{cancel}
\renewcommand{\CancelColor}{\color{coulExo}}
\newcommand{\fois}[1]{\times{\color{coulExo}#1}}
\newcommand{\barre}[1]{\cancel{#1}}
\newlength{\lycee@eqh}\newlength{\lycee@eqd}
\newsavebox{\lycee@eqa}\newsavebox{\lycee@eqb}\newsavebox{\lycee@eqc}
\newcommand{\lycee@eqcol}[1]{\hbox{$\renewcommand{\arraystretch}{1.35}%
  \begin{array}[t]{@{}>{\displaystyle}r@{}>{\displaystyle{}}l@{}}#1\end{array}$}}
\newcommand{\lycee@eqmesure}[1]{%
  \ifdim\ht#1>\lycee@eqh\setlength{\lycee@eqh}{\ht#1}\fi
  \ifdim\dp#1>\lycee@eqd\setlength{\lycee@eqd}{\dp#1}\fi}
\newcommand{\lycee@eqtrait}{\hspace{1.2em}\tikz[baseline=0pt]\draw[dash pattern=on 2pt off 2pt,
  line width=0.5pt] (0,\lycee@eqh) -- (0,-\lycee@eqd);\hspace{1.2em}}
\newcommand{\eqproduit}[2]{\par\smallskip\noindent C'est une équation produit :
  $A\times B=0 \iff A=0$ ou $B=0$.\par\nopagebreak\smallskip
  \begingroup\setlength{\lycee@eqh}{0pt}\setlength{\lycee@eqd}{0pt}%
  \setbox\lycee@eqa=\lycee@eqcol{#1}\setbox\lycee@eqb=\lycee@eqcol{#2}%
  \lycee@eqmesure\lycee@eqa\lycee@eqmesure\lycee@eqb%
  \noindent\hspace*{\retraitcalculs}\usebox\lycee@eqa\lycee@eqtrait\usebox\lycee@eqb\endgroup\par\smallskip}
\newcommand{\eqproduittrois}[3]{\par\smallskip\noindent C'est une équation produit :
  $A\times B\times C=0 \iff A=0$ ou $B=0$ ou $C=0$.\par\nopagebreak\smallskip
  \begingroup\setlength{\lycee@eqh}{0pt}\setlength{\lycee@eqd}{0pt}%
  \setbox\lycee@eqa=\lycee@eqcol{#1}\setbox\lycee@eqb=\lycee@eqcol{#2}\setbox\lycee@eqc=\lycee@eqcol{#3}%
  \lycee@eqmesure\lycee@eqa\lycee@eqmesure\lycee@eqb\lycee@eqmesure\lycee@eqc%
  \noindent\hspace*{\retraitcalculs}\usebox\lycee@eqa\lycee@eqtrait\usebox\lycee@eqb\lycee@eqtrait\usebox\lycee@eqc\endgroup
  \par\smallskip}

\tikzset{grille/.style={coulGrille,line width=0.4pt}}
\newcommand{\quadrillage}[4]{\draw[grille,step=1] (#1,#2) grid (#3,#4);}
\tikzset{
  croix/.style={cross out,draw,line width=0.6pt,inner sep=0pt,minimum size=4pt},
  croixdroite/.style={croix,rotate=45,line width=0.8pt},
}
\newcommand{\pt}[3]{\path (#1) node[croix]{} node[#2,font=\small,inner sep=2.5pt]{$#3$};}

\newcommand{\lycee@repere}[4]{%
  \draw[grille,step=1] (#1,#2) grid (#3,#4);
  \draw[-{Stealth[length=2mm]},thick] (#1,0)--({#3+0.4},0) node[below left=-1pt]{$x$};
  \draw[-{Stealth[length=2mm]},thick] (0,#2)--(0,{#4+0.4}) node[below left=-1pt]{$y$};
  \node[below left,font=\scriptsize,inner sep=1.5pt] at (0,0) {$0$};}
\newcommand{\lycee@gradx}[1]{\foreach \k/\l in {#1}{%
  \draw (\k,0.07)--(\k,-0.07) node[below,font=\scriptsize,inner sep=1.5pt]{$\l$};}}
\newcommand{\lycee@grady}[1]{\foreach \k/\l in {#1}{%
  \draw (0.07,\k)--(-0.07,\k) node[left,font=\scriptsize,inner sep=1.5pt]{$\l$};}}
\newcommand{\lycee@intff}[2]{\left[#1\,;#2\right]}
\newcommand{\lycee@intfo}[2]{\left[#1\,;#2\right[}
\newcommand{\lycee@intof}[2]{\left]#1\,;#2\right]}
\newcommand{\lycee@intoo}[2]{\left]#1\,;#2\right[}
\newcommand{\lycee@ens}[1]{\left\{#1\right\}}
\AtBeginDocument{%
  \providecommand{\repere}{\lycee@repere}%
  \providecommand{\gradx}{\lycee@gradx}%
  \providecommand{\grady}{\lycee@grady}%
  \providecommand{\Cf}{\mathcal{C}\sb{\scriptscriptstyle f}}%
  \providecommand{\Cg}{\mathcal{C}\sb{\scriptscriptstyle g}}%
  \providecommand{\intff}{\lycee@intff}%
  \providecommand{\intfo}{\lycee@intfo}%
  \providecommand{\intof}{\lycee@intof}%
  \providecommand{\intoo}{\lycee@intoo}%
  \providecommand{\ens}{\lycee@ens}}

\newlist{questions}{enumerate}{3}
\setlist[questions,1]{label=\arabic*\textdegree),leftmargin=*,itemsep=2pt,topsep=3pt}
\setlist[questions,2]{label=\alph*),leftmargin=*,itemsep=1pt,topsep=2pt}
\setlist[questions,3]{label=\roman*),leftmargin=*,itemsep=1pt,topsep=1pt}
\setlist[itemize]{label=\textbullet,leftmargin=*,itemsep=2pt,topsep=3pt}

\newcommand{\besoinplace}[1]{\par\penalty\@M\begingroup
  \dimen@=#1\relax\dimen@ii=\pagegoal\advance\dimen@ii-\pagetotal
  \ifdim\dimen@>\dimen@ii\ifdim\dimen@ii>\z@\vfil\fi\break\fi\endgroup}

\newcommand{\lycee@pied}{\thepage/\pageref{LastPage}}
\newcommand{\entete}[2]{%
  \pagestyle{fancy}\fancyhf{}%
  \fancyhead[L]{\footnotesize\color{coulPart}\textbf{#1}}%
  \fancyhead[R]{\footnotesize\color{coulPart}\textbf{#2}}%
  \fancyfoot[C]{\footnotesize\lycee@pied}%
  \renewcommand{\headrulewidth}{0.4pt}%
  \renewcommand{\headrule}{\hbox to\headwidth{%
    \color{coulPart!40}\leaders\hrule height \headrulewidth\hfill}}}

\RequirePackage{ccicons}
\newcommand{\lycee@licence}{{\scriptsize\color{gray}%
  \copyright\ Luc Destombe --- \ccbyncsa\ CC BY-NC-SA 4.0}}
\newif\iflycee@licence \lycee@licencetrue
\newcommand{\sanslicence}{\lycee@licencefalse}
\AtBeginDocument{\iflycee@licence\AddToHookNext{shipout/foreground}{%
  \put(0,0){\rlap{\kern\dimexpr1in+\hoffset+\oddsidemargin+\textwidth\relax
    \llap{\raisebox{-\dimexpr1in+\voffset+\topmargin+\headheight+\headsep
      +\textheight+\footskip\relax}{\lycee@licence}}}}}\fi}

\newcommand{\lycee@site}{faire-des-maths.fr}
\newcommand{\lycee@version}{{\scriptsize\color{gray}\lycee@site
  \ifdefined\lyceeversion\ --- Version~\lyceeversion\fi}}
\AtBeginDocument{\iflycee@licence\AddToHookNext{shipout/foreground}{%
  \put(0,0){\rlap{\kern\dimexpr1in+\hoffset+\oddsidemargin\relax
    \raisebox{-\dimexpr1in+\voffset+\topmargin+\headheight+\headsep
      +\textheight+\footskip\relax}{\lycee@version}}}}\fi}

\newcommand{\formulesserrees}{%
  \setlength{\abovedisplayskip}{3pt}\setlength{\belowdisplayskip}{3pt}%
  \setlength{\abovedisplayshortskip}{2pt}\setlength{\belowdisplayshortskip}{3pt}}

\setlength{\parindent}{0pt}

\RequirePackage[implicit=false,hidelinks,pdfencoding=auto,psdextra,bookmarksopen,
  bookmarksopenlevel=1,pdfcreator={},pdfproducer={}]{hyperref}
\RequirePackage{bookmark}
\hypersetup{pdfauthor={Luc Destombe},pdfkeywords={CC BY-NC-SA 4.0}}
\newcounter{lycee@signet}
\newcommand{\signet}[3][]{\stepcounter{lycee@signet}%
  \edef\lycee@dest{\if\relax\detokenize{#1}\relax
    signet.\arabic{lycee@signet}\else#1\fi}%
  \hypertarget{\lycee@dest}{}\bookmark[level=#2,dest=\lycee@dest]{#3}}
\pdfstringdefDisableCommands{%
  \def\R{R}\def\vect#1{#1}\def\Cf{Cf}\def\Cg{Cg}\def\fois#1{×#1}%
  \def\barre#1{#1}\def\intervalle#1#2{[#1;#2]}\def\intff#1#2{[#1;#2]}%
  \def\textsuperscript#1{#1}\def\dfrac#1#2{#1/#2}\def\frac#1#2{#1/#2}%
  \def\vec#1{#1⃗}}%
\RequirePackage[official]{eurosym}

\geometry{top=0.8cm,bottom=1.3cm,left=1cm,right=1cm,footskip=0.6cm}

\pagestyle{fancy}
\fancyhf{}
\renewcommand{\headrulewidth}{0pt}
\fancyfoot[C]{\footnotesize Page \thepage{} / \pageref{LastPage}}

\newcommand{\titrefiche}[2]{%
  \begin{center}
    \Large\textbf{#1}\\[2pt]
    \large\textbf{#2}\\[2pt]
  \end{center}
  \vspace{0.10cm}}

\setlength{\columnseprule}{0.4pt}
\setlength{\columnsep}{15pt}
\renewcommand{\columnseprulecolor}{\color{black!45}}
\setlength{\multicolsep}{2pt}

\newcounter{partiefiche}
\renewcommand{\thepartiefiche}{\Roman{partiefiche}}
\newif\ifapresbandeau
\newcommand{\lycee@bandeau}[1]{%
  \par\penalty-600
  \vspace{0.08cm}%
  \noindent\colorbox{coulPart}{%
    \parbox{\dimexpr\linewidth-2\fboxsep\relax}{%
      \color{white}\bfseries\raggedright\strut#1\strut}}%
  \nopagebreak\par\nopagebreak
  \vspace{0.06cm}\nopagebreak
  \apresbandeautrue}
\newcommand{\partiefiche}[1]{%
  \refstepcounter{partiefiche}%
  \lycee@bandeau{\signet{1}{\thepartiefiche{} --- #1}\thepartiefiche{} --- #1}}
\newcommand{\partiefichesansnum}[1]{\lycee@bandeau{\signet{1}{#1}#1}}

\newcounter{exo}
\newlength{\espaceexo}\setlength{\espaceexo}{7pt}
\newlength{\espacetitre}\setlength{\espacetitre}{2pt}
\newcommand{\exo}[1][\relax]{%
  \par
  \ifapresbandeau\apresbandeaufalse\else\penalty-150\addvspace{\espaceexo}\fi
  \refstepcounter{exo}%
  \noindent\signet[exercice-\arabic{exo}]{2}{Exercice \arabic{exo}}%
  \textbf{\color{coulExo}\underline{Exercice \theexo}}%
  \ifx\relax#1\relax\else\textbf{ : #1}\fi
  \par\nopagebreak\vspace{\espacetitre}\nopagebreak
  \ignorespaces}

\setlist[enumerate]{nosep,leftmargin=*,itemsep=2pt,topsep=2pt}
\setlist[itemize]{nosep,leftmargin=*,topsep=2pt}
\newcommand{\choix}[3]{\par\hspace*{1em}%
  \makebox[0.3\linewidth][l]{a) #1}\makebox[0.3\linewidth][l]{b) #2}c) #3}
\newcommand{\deux}[2]{\par\hspace*{0.6em}%
  \makebox[0.48\linewidth][l]{#1}#2}
\newcommand{\trois}[3]{\par\hspace*{0.6em}%
  \makebox[0.32\linewidth][l]{#1}\makebox[0.32\linewidth][l]{#2}#3}

  \renewcommand{\exo}[1][\relax]{%
    \ifapresbandeau\apresbandeaufalse\else\par\penalty-150\fi
    \refstepcounter{exo}%
    \vspace{0.05cm}%
    \noindent\signet[exercice-\arabic{exo}]{2}{Exercice \arabic{exo}}%
    \textbf{\color{coulExo}\underline{Exercice \theexo}}%
    \ifx\relax#1\relax\else\textbf{ : #1}\fi%
    \par\nopagebreak
    \ignorespaces}
  \newcommand{\rappel}[1]{%
    \par\vspace{0.06cm}\nopagebreak
    \noindent\fcolorbox{black!35}{black!5}{%
      \parbox{\dimexpr\linewidth-2\fboxsep-2\fboxrule\relax}{%
        \small\textbf{Rappel.} #1}}%
    \par\vspace{0.06cm}\nopagebreak}
  \setlist[enumerate]{nosep,leftmargin=*}
  \setlist[itemize]{nosep,leftmargin=*}
  \setlength{\multicolsep}{3pt plus 1pt minus 1pt}
  \setlength{\premulticols}{10pt}
  \setlength{\postmulticols}{10pt}
  \newlist{expr}{enumerate}{1}
  \setlist[expr]{label={\Alph*\ $=$},nosep,leftmargin=*,itemsep=0pt}

\renewenvironment{center}{\par\addvspace{3pt}\centering}{\par\addvspace{3pt}}
\formulesserrees
\AtBeginDocument{\formulesserrees}
\clubpenalty=10000
\widowpenalty=10000
\displaywidowpenalty=10000
\brokenpenalty=10000
\setlength{\parskip}{0pt}

\RequirePackage{ragged2e}
\setlength{\RaggedRightRightskip}{0pt plus 1fil}
\AtBeginDocument{\RaggedRight\binoppenalty=10000 \relpenalty=10000 }
\g@addto@macro\@parboxrestore{\RaggedRight}
\makeatother

\newcommand{\coord}[2]{\begin{pmatrix}#1\\#2\end{pmatrix}}
\ShorthandsOff
\tikzset{fig/.style={line width=0.8pt,black!75}}
\ShorthandsOn

\begin{document}

\titrefiche{1\textsuperscript{re} spécialité --- Géométrie repérée et produit scalaire}{Fiche d'exercices du chapitre}

{\small\itshape À partir de la partie IV, le plan est muni d'un repère orthonormé.\par}

\begin{multicols*}{2}\raggedcolumns

\partiefiche{Coordonnées d'un vecteur, somme de deux vecteurs}

\exo[Coordonnées de vecteurs]
\begin{enumerate}
  \item Donner les coordonnées des vecteurs $\vec{u}=-2\vec{\imath}+7\vec{\jmath}$ et
        $\vec{v}=5\vec{\jmath}$.
  \item Pour chaque couple de points, calculer les coordonnées de $\vect{AB}$, puis
        celles de $\vect{BA}$ : a)~$A(2\,;5)$ et $B(6\,;1)$ ; b)~$A(-3\,;2)$ et
        $B(1\,;-4)$ ; c)~$A(-1\,;-5)$ et $B(-7\,;-2)$ ;
        d)~$A\left(\frac{1}{3}\,;2\right)$ et $B\left(-1\,;\frac{5}{3}\right)$.
\end{enumerate}

\exo[Calculer avec les coordonnées]
On considère les vecteurs $\vec{u}\coord{2}{-5}$, $\vec{v}\coord{-3}{1}$ et
$\vec{w}\coord{4}{0}$.
\begin{enumerate}
  \item Calculer les coordonnées de $\vec{u}+\vec{v}$, de $\vec{u}-\vec{v}$ et de
        $\vec{v}-\vec{u}$.
  \item Calculer les coordonnées du vecteur $\vec{t}$ tel que $\vec{u}+\vec{t}=\vec{w}$.
\end{enumerate}

\exo[Relation de Chasles]
On considère les points $A(0\,;3)$, $B(4\,;1)$ et $C(-2\,;-1)$.
\begin{enumerate}
  \item Calculer les coordonnées de $\vect{AB}+\vect{BC}$ et celles de $\vect{AC}$.
        Que constate-t-on ?
  \item Calculer les coordonnées de $\vect{AB}-\vect{AC}$. Quel vecteur, nommé avec
        deux des points $A$, $B$, $C$, retrouve-t-on ?
\end{enumerate}

\exo[Trajet d'un drone]
Un drone décolle du point $O(0\,;0)$ (unité : la dizaine de mètres). Il se déplace
successivement selon les vecteurs $\vec{u}\coord{3}{2}$, $\vec{v}\coord{-1}{4}$ et
$\vec{w}\coord{2}{-3}$.
\begin{enumerate}
  \item Calculer les coordonnées des points où le drone s'arrête après chaque
        déplacement.
  \item Calculer les coordonnées du vecteur $\vec{u}+\vec{v}+\vec{w}$. Que représente-t-il ?
  \item Quel déplacement ramène le drone à son point de départ ?
\end{enumerate}

\partiefiche{Parallélogramme et milieu d'un segment}

\exo[Point défini par une égalité de vecteurs]
On considère les points $A(1\,;-2)$, $B(4\,;1)$, $C(2\,;5)$ et le vecteur
$\vec{u}\coord{-3}{4}$. Calculer les coordonnées :
\begin{enumerate}[label=\alph*)]
  \item du point $E$ tel que $\vect{AE}=\vec{u}$ ;
  \item du point $D$ tel que $ABCD$ soit un parallélogramme ;
  \item du point $F$ tel que $\vect{FA}=\vect{BC}$.
\end{enumerate}

\exo[Milieux]
\begin{enumerate}
  \item On considère les points $R(-4\,;3)$ et $S(2\,;7)$. Calculer les coordonnées du
        milieu $I$ de $[RS]$.
  \item Le point $K(1\,;-2)$ est le milieu du segment $[RT]$. Calculer les coordonnées
        de $T$.
  \item Calculer les coordonnées du milieu de $[EF]$, avec $E\left(\frac{1}{3}\,;2\right)$
        et $F(1\,;-1)$.
\end{enumerate}

\exo[Parallélogramme ou non ?]
Pour chaque quadrilatère, dire s'il est un parallélogramme, en comparant les milieux
de ses diagonales : a)~$KLMN$ avec $K(-2\,;-1)$, $L(3\,;0)$, $M(5\,;4)$ et $N(0\,;3)$ ;
b)~$KLMV$ avec $V(1\,;3)$.

\exo[Quatrième sommet]
On considère les points $A(-1\,;1)$, $B(3\,;-1)$ et $C(4\,;2)$.
\begin{enumerate}
  \item Calculer les coordonnées du point $D$ tel que $ABCD$ soit un parallélogramme.
  \item Calculer les coordonnées du centre $O$ de ce parallélogramme.
  \item Calculer les coordonnées de $\vect{AB}+\vect{AD}$ et celles de $\vect{AC}$.
        Quelle règle retrouve-t-on ?
  \item Calculer les coordonnées du point $E$ tel que $B$ soit le milieu de $[AE]$.
\end{enumerate}

\partiefiche{Produit d'un vecteur par un nombre réel, colinéarité}

\exo[Vecteurs colinéaires ou non]
\begin{enumerate}
  \item Pour chaque couple de vecteurs, dire s'ils sont colinéaires :
        a)~$\vec{u}\coord{3}{-4}$ et $\vec{v}\coord{-9}{12}$ ; b)~$\vec{u}\coord{5}{2}$ et
        $\vec{w}\coord{7}{3}$ ; c)~$\vec{a}\coord{\sqrt{2}}{2}$ et $\vec{b}\coord{1}{\sqrt{2}}$.
  \item Calculer les nombres réels $m$ pour lesquels $\vec{r}\coord{m}{3}$ et
        $\vec{s}\coord{27}{m}$ sont colinéaires.
\end{enumerate}

\exo[Combinaisons de vecteurs]
On considère les vecteurs $\vec{u}\coord{2}{-1}$ et $\vec{v}\coord{-3}{4}$.
\begin{enumerate}
  \item Calculer les coordonnées de $3\vec{u}-2\vec{v}$ et de $\frac{1}{2}\vec{u}+\vec{v}$.
  \item Calculer les coordonnées du vecteur $\vec{w}$ tel que $2\vec{w}+\vec{u}=\vec{v}$.
\end{enumerate}

\exo[Parallélisme et alignement]
On considère les points $A(-2\,;1)$, $B(2\,;3)$, $C(0\,;-2)$, $D(8\,;2)$ et $E(8\,;6)$.
\begin{enumerate}
  \item Les droites $(AB)$ et $(CD)$ sont-elles parallèles ?
  \item Les points $A$, $B$ et $E$ sont-ils alignés ? Et les points $A$, $B$ et $C$ ?
  \item Calculer l'ordonnée $y$ du point $M(10\,;y)$ de la droite $(AB)$.
\end{enumerate}

\exo[Droite des milieux]
On considère les points $A(-2\,;0)$, $B(4\,;2)$ et $C(2\,;6)$ ; $I$ est le milieu de
$[AC]$ et $J$ celui de $[BC]$.
\begin{enumerate}
  \item Calculer les coordonnées de $I$ et de $J$.
  \item Démontrer que $\vect{IJ}=\frac{1}{2}\vect{AB}$.
  \item En déduire que les droites $(IJ)$ et $(AB)$ sont parallèles.
  \item Le point $E(-5\,;-1)$ est-il sur la droite $(AB)$ ?
\end{enumerate}

\partiefiche{Norme d'un vecteur, équation d'un cercle}

\exo[Normes et distances]
\begin{enumerate}
  \item Calculer la norme du vecteur $\vec{u}\coord{-7}{24}$.
  \item On considère les points $C(1\,;-2)$ et $D(-3\,;6)$. Calculer $CD$ et l'écrire
        sous la forme $a\sqrt{b}$, avec deux nombres entiers $a$ et $b$.
  \item Pour chaque triangle, donner sa nature : a)~$EFG$ avec $E(1\,;-1)$, $F(3\,;3)$
        et $G(-1\,;5)$ ; b)~$KLM$ avec $K(0\,;1)$, $L(4\,;4)$ et $M(-3\,;-3)$.
\end{enumerate}

\exo[Équations de cercles]
\begin{enumerate}
  \item Pour chaque cercle, lire le centre et le rayon :
        a)~$(x-4)^{2}+(y+1)^{2}=9$ ; b)~$x^{2}+(y-2)^{2}=5$.
  \item Écrire une équation du cercle $\mathcal{C}_{1}$ de centre $\Omega(-3\,;2)$ et de
        rayon $5$.
  \item Écrire une équation du cercle $\mathcal{C}_{2}$ de centre $A(2\,;1)$ qui passe
        par $B(-1\,;5)$. Les points $C(5\,;5)$ et $D(-2\,;3)$ sont-ils sur $\mathcal{C}_{2}$ ?
  \item Écrire une équation du cercle $\mathcal{C}_{3}$ de diamètre $[EF]$, avec
        $E(-4\,;3)$ et $F(2\,;-5)$.
\end{enumerate}

\exo[Reconnaître un cercle]
Pour chaque équation, compléter les carrés (forme canonique) pour dire si c'est une
équation de cercle et, si oui, donner son centre et son rayon :
a)~$x^{2}+y^{2}-6x+4y-3=0$ ; b)~$x^{2}+y^{2}+8x-2y+8=0$ ;
c)~$x^{2}+y^{2}-2x+6y+11=0$.
\par
Calculer ensuite les coordonnées du milieu du segment qui joint les centres des
cercles a) et b).

\exo[Zone couverte par une antenne]
Une antenne relais placée en $\Omega(2\,;3)$ couvre un disque de rayon $5$ (unité : le
kilomètre).
\begin{enumerate}
  \item Écrire une équation du cercle $\mathcal{C}$ qui borde la zone couverte.
  \item Les villages $V(5\,;7)$ et $W(-2\,;-1)$ sont-ils couverts ?
  \item Calculer les coordonnées du point $V'$ tel que $\Omega$ soit le milieu de $[VV']$.
  \item Une route rectiligne a pour équation $y=-1$. Calculer les coordonnées des
        points de la route situés sur $\mathcal{C}$. Retrouve-t-on $V'$ ?
\end{enumerate}

\partiefiche{Définition du produit scalaire}

\exo[Produit scalaire avec un angle]
\begin{enumerate}[label=\alph*)]
  \item $\norme{\vec{u}}=4$, $\norme{\vec{v}}=5$ et l'angle de $\vec{u}$ et $\vec{v}$ mesure
        $60^{\circ}$. Calculer $\vec{u}\cdot\vec{v}$.
  \item Dans le triangle $ABC$, $AB=6$, $AC=2$ et $\widehat{BAC}=150^{\circ}$. Calculer
        $\vect{AB}\cdot\vect{AC}$.
  \item Dans le triangle $DEF$, $DE=3$, $DF=4\sqrt{2}$ et $\widehat{EDF}=45^{\circ}$.
        Calculer $\vect{DE}\cdot\vect{DF}$.
  \item $KLM$ est un triangle équilatéral de côté $8$. Calculer $\vect{KL}\cdot\vect{KM}$,
        puis $\vect{KL}\cdot\vect{LM}$.
\end{enumerate}

\exo[Vrai ou faux]
Pour chaque phrase, dire si elle est vraie ou fausse, en justifiant par une phrase.
\begin{enumerate}[label=\alph*)]
  \item Si $\vec{u}\cdot\vec{v}<0$, alors l'angle de $\vec{u}$ et $\vec{v}$ est obtus.
  \item Si $\vec{u}\cdot\vec{v}=0$, alors $\vec{u}$ ou $\vec{v}$ est le vecteur nul.
  \item Pour tout vecteur $\vec{u}$, $\vec{u}\cdot\vec{u}=\norme{\vec{u}}^{2}$.
  \item Deux vecteurs de norme $3$ ont toujours un produit scalaire égal à $9$.
\end{enumerate}

\exo[Hexagone régulier]
$ABCDEF$ est un hexagone régulier de centre $O$ et de côté $2$ : les triangles $OAB$,
$OBC$, \dots{} sont équilatéraux. Calculer $\vect{OA}\cdot\vect{OB}$,
$\vect{OA}\cdot\vect{OC}$, $\vect{OA}\cdot\vect{OD}$ et $\vect{AB}\cdot\vect{AF}$.

\exo[Travail d'une force]
En physique, le \textbf{travail} d'une force constante $\vec{F}$ le long d'un
déplacement rectiligne $\vect{AB}$ est le produit scalaire
$W=\vec{F}\cdot\vect{AB}$ (en joules, avec $F$ en newtons et $AB$ en mètres). Calculer
le travail dans chaque situation.
\begin{enumerate}[label=\alph*)]
  \item Un chariot est poussé sur $15$~m par une force de $40$~N, inclinée de
        $30^{\circ}$ par rapport au sol (valeur exacte, puis arrondie à l'unité).
  \item Un téléski tire un skieur sur $300$~m avec une force de $200$~N ; la perche fait
        un angle de $20^{\circ}$ avec la piste (arrondi au joule).
  \item Une valise est roulée sur $50$~m de sol horizontal ; son poids, de $100$~N, est
        vertical.
  \item Un frottement de $30$~N, de sens contraire au déplacement, freine un palet sur
        $15$~m.
\end{enumerate}

\partiefiche{Produit scalaire et projeté orthogonal}

\exo[Dans un rectangle]
\begin{minipage}[t]{0.6\linewidth}\raggedright
$ABCD$ est un rectangle avec $AB=6$ et $AD=4$ ; $I$ est le milieu de $[CD]$ et $J$
celui de $[BC]$. Calculer :
\end{minipage}\hfill
\begin{minipage}[t]{0.36\linewidth}
\vspace{0pt}\centering
\begin{tikzpicture}[scale=0.32]
  \draw[fig] (0,0) -- (6,0) -- (6,4) -- (0,4) -- cycle;
  \pt{0,0}{below left}{A} \pt{6,0}{below right}{B} \pt{6,4}{above right}{C}
  \pt{0,4}{above left}{D} \pt{3,4}{above}{I} \pt{6,2}{right}{J}
\end{tikzpicture}
\end{minipage}
\par
a)~$\vect{AB}\cdot\vect{AC}$ ; b)~$\vect{AB}\cdot\vect{AD}$ ; c)~$\vect{AB}\cdot\vect{CD}$ ;
d)~$\vect{AB}\cdot\vect{AI}$ ; e)~$\vect{AD}\cdot\vect{AJ}$.
\par
f)~Dans un triangle $EFG$, $H$ est le projeté orthogonal de $G$ sur $(EF)$, le point
$E$ est entre $H$ et $F$, $EF=7$ et $EH=2$. Calculer $\vect{EF}\cdot\vect{EG}$.

\exo[Sur une droite graduée]
Sur une droite graduée, les points $A$, $B$, $C$, $D$ et $E$ ont pour abscisses $-2$,
$0$, $1$, $3$ et $4$. Calculer $\vect{AB}\cdot\vect{AD}$, $\vect{BC}\cdot\vect{BE}$,
$\vect{CA}\cdot\vect{CD}$ et $\vect{AE}\cdot\vect{DB}$.

\exo[Dans un carré]
$ABCD$ est un carré de côté $4$ ; $I$ est le milieu de $[AB]$ et $J$ celui de $[BC]$.
\begin{enumerate}
  \item Calculer $\vect{AB}\cdot\vect{AJ}$ et $\vect{AI}\cdot\vect{AJ}$.
  \item Calculer $\vect{DA}\cdot\vect{DJ}$ (projeter $J$ sur la droite $(DA)$).
  \item Démontrer, avec la relation de Chasles, que $\vect{IJ}=\frac{1}{2}\vect{AC}$.
\end{enumerate}

\partiefiche{Règles de calcul}

\exo[Calculer avec les règles]
On sait que $\norme{\vec{u}}=2$, $\norme{\vec{v}}=5$ et $\vec{u}\cdot\vec{v}=3$. Calculer :
a)~$(4\vec{u})\cdot(-\vec{v})$ ; b)~$\vec{u}\cdot(\vec{u}-\vec{v})$ ;
c)~$\norme{\vec{u}+\vec{v}}$ ; d)~$\norme{\vec{u}-\vec{v}}$ ;
e)~$(\vec{u}+2\vec{v})\cdot(\vec{u}-\vec{v})$.

\exo[Diagonales d'un losange]
\begin{enumerate}
  \item Démontrer que, pour tous vecteurs $\vec{u}$ et $\vec{v}$,
        $(\vec{u}+\vec{v})\cdot(\vec{u}-\vec{v})=\norme{\vec{u}}^{2}-\norme{\vec{v}}^{2}$.
  \item $ABCD$ est un losange. Écrire $\vect{AC}$ et $\vect{DB}$ avec $\vect{AB}$ et
        $\vect{AD}$.
  \item En déduire que les diagonales d'un losange sont perpendiculaires.
\end{enumerate}

\exo[Angle obtus dans un rectangle]
$ABCD$ est un rectangle avec $AB=6$ et $AD=4$. Le point $K$ est sur $[AB]$, avec
$AK=4$, et le point $L$ sur $[BC]$, avec $BL=1$.
\begin{enumerate}
  \item Écrire $\vect{KD}$ avec $\vect{KA}$ et $\vect{AD}$, puis $\vect{KL}$ avec
        $\vect{KB}$ et $\vect{BL}$.
  \item Expliquer pourquoi $\vect{KA}\cdot\vect{BL}=0$ et $\vect{AD}\cdot\vect{KB}=0$.
  \item Calculer $\vect{KA}\cdot\vect{KB}$ et $\vect{AD}\cdot\vect{BL}$.
  \item En déduire $\vect{KD}\cdot\vect{KL}$. L'angle $\widehat{DKL}$ est-il aigu ou obtus ?
\end{enumerate}

\partiefiche{Produit scalaire dans un repère orthonormé}

\exo[Produit scalaire et orthogonalité]
\begin{enumerate}[label=\alph*)]
  \item Calculer $\vec{u}\cdot\vec{v}$ pour $\vec{u}\coord{5}{-2}$ et $\vec{v}\coord{3}{4}$,
        puis pour $\vec{u}\,{}'\coord{\frac{1}{3}}{2}$ et $\vec{v}\,{}'\coord{6}{-\frac{1}{2}}$.
  \item Les vecteurs $\vec{a}\coord{-4}{6}$ et $\vec{b}\coord{3}{2}$ sont-ils orthogonaux ?
  \item On considère les points $A(2\,;-1)$, $B(4\,;3)$ et $C(-2\,;1)$. Le triangle $ABC$
        est-il rectangle en $A$ ?
  \item Calculer le nombre réel $m$ pour lequel $\vec{r}\coord{m}{2}$ et $\vec{s}\coord{4}{-6}$
        sont orthogonaux.
\end{enumerate}

\exo[Rectangle en quel sommet ?]
Pour chaque triangle, chercher s'il est rectangle, et en quel sommet ; conclure par
une phrase. a)~$KLM$ avec $K(1\,;1)$, $L(5\,;3)$ et $M(0\,;3)$ ; b)~$RST$ avec
$R(-2\,;0)$, $S(2\,;2)$ et $T(4\,;-2)$.

\exo[Vecteurs de la base]
\begin{enumerate}
  \item Calculer $\vec{u}\cdot\vec{\imath}$ et $\vec{u}\cdot\vec{\jmath}$ pour
        $\vec{u}\coord{-2}{7}$.
  \item Un vecteur $\vec{w}$ vérifie $\vec{w}\cdot\vec{\imath}=4$ et
        $\vec{w}\cdot\vec{\jmath}=-1$. Donner ses coordonnées.
  \item Vérifier que $\vec{v}\coord{7}{2}$ est orthogonal à $\vec{u}$.
  \item Le vecteur $\vec{t}\coord{14}{4}$ est-il orthogonal à $\vec{u}$ ? Est-il
        colinéaire à $\vec{v}$ ?
\end{enumerate}

\exo[Où placer le point ?]
On considère les points $A(1\,;2)$ et $B(5\,;0)$, et un point $M(x\,;3)$ de la droite
d'équation $y=3$.
\begin{enumerate}
  \item Calculer la distance $AB$.
  \item Calculer $x$ pour que le triangle $ABM$ soit rectangle en $A$.
  \item Démontrer que le triangle $ABM$ est rectangle en $M$ quand $x^{2}-6x+8=0$.
  \item En déduire les points $M$ de la droite pour lesquels $ABM$ est rectangle en $M$.
\end{enumerate}

\partiefiche{Produit scalaire et normes}

\exo[Produit scalaire avec des normes]
\begin{enumerate}[label=\alph*)]
  \item $\norme{\vec{u}}=4$, $\norme{\vec{v}}=6$ et $\norme{\vec{u}+\vec{v}}=8$. Calculer
        $\vec{u}\cdot\vec{v}$.
  \item Dans le triangle $ABC$, $AB=5$, $AC=8$ et $BC=7$. Calculer
        $\vect{AB}\cdot\vect{AC}$, $\vect{BA}\cdot\vect{BC}$ et $\vect{CA}\cdot\vect{CB}$.
  \item $\norme{\vec{u}}=3$, $\norme{\vec{v}}=2$ et $\norme{\vec{u}-\vec{v}}=4$. Calculer
        $\vec{u}\cdot\vec{v}$ (avec la deuxième identité remarquable).
\end{enumerate}

\exo[Diagonales d'un parallélogramme]
$ABCD$ est un parallélogramme avec $AB=6$, $AD=4$ et $BD=5$.
\begin{enumerate}
  \item Calculer $\vect{AB}\cdot\vect{AD}$ dans le triangle $ABD$.
  \item Calculer $AC^{2}$, avec $\vect{AC}=\vect{AB}+\vect{AD}$, puis $AC$.
  \item Vérifier que $AC^{2}+BD^{2}=2\left(AB^{2}+AD^{2}\right)$.
\end{enumerate}

\exo[Une voile triangulaire]
Une voile triangulaire $ABC$ a pour côtés $AB=3$~m, $AC=4$~m et $BC=6$~m.
\begin{enumerate}
  \item Calculer $\vect{AB}\cdot\vect{AC}$, $\vect{BA}\cdot\vect{BC}$ et
        $\vect{CA}\cdot\vect{CB}$.
  \item En déduire la nature (aigu ou obtus) de chaque angle de la voile, en une phrase.
\end{enumerate}

\partiefiche{Vecteur normal et équation cartésienne d'une droite}

\exo[Équation avec un vecteur normal]
\begin{enumerate}[label=\alph*)]
  \item Pour chaque droite, lire un vecteur normal et un vecteur directeur :
        $(d_{1})$ : $5x-3y+2=0$ ; $(d_{2})$ : $-2x+y=0$ ; $(d_{3})$ : $y=4x-1$.
  \item Trouver une équation cartésienne de la droite qui passe par $A(1\,;-3)$, de
        vecteur normal $\vec{n}\coord{2}{5}$.
  \item On considère les points $B(2\,;1)$, $C(-1\,;3)$ et $D(3\,;-2)$. Trouver une
        équation de la hauteur issue de $B$ dans le triangle $BCD$.
  \item Trouver une équation de la médiatrice de $[EF]$, avec $E(-3\,;2)$ et $F(1\,;6)$.
\end{enumerate}

\exo[Parallèles ou perpendiculaires ?]
Pour chaque couple de droites, dire si elles sont parallèles, perpendiculaires, ou ni
l'un ni l'autre : a)~$3x+4y-1=0$ et $4x-3y+2=0$ ; b)~$3x+4y-1=0$ et $-6x-8y+5=0$ ;
c)~$x-2y+3=0$ et $y=3x+1$ ; d)~$x=4$ et $y=-2$.

\exo[Tangentes à un cercle]
On considère le cercle $\mathcal{C}$ de centre $\Omega(1\,;-2)$ et de rayon $5$. La
tangente à $\mathcal{C}$ en un point $A$ est la droite perpendiculaire à $(\Omega A)$
en $A$.
\begin{enumerate}
  \item Écrire une équation de $\mathcal{C}$ et vérifier que $A(4\,;2)$ et $B(-2\,;2)$
        sont sur $\mathcal{C}$.
  \item Trouver une équation de la tangente à $\mathcal{C}$ en $A$, puis en $B$.
  \item Ces deux tangentes sont-elles perpendiculaires ?
\end{enumerate}

\exo[Les trois hauteurs]
On considère le triangle $ABC$ avec $A(0\,;4)$, $B(-2\,;0)$ et $C(4\,;0)$.
\begin{enumerate}
  \item Trouver une équation de la hauteur issue de $A$, puis de la hauteur issue de $B$.
  \item Calculer les coordonnées du point $H$ commun à ces deux hauteurs.
  \item Démontrer que les droites $(CH)$ et $(AB)$ sont perpendiculaires. Que
        peut-on en conclure pour les trois hauteurs du triangle ?
  \item Calculer $\vect{AB}\cdot\vect{AC}$. L'angle $\widehat{BAC}$ est-il aigu ?
\end{enumerate}

\partiefiche{Projeté orthogonal d'un point sur une droite}

\exo[Projeté orthogonal]
\begin{enumerate}
  \item On considère la droite $(d)$ d'équation $x+2y-5=0$ et le point $M(3\,;6)$.
        a)~Trouver une équation de la perpendiculaire $(\Delta)$ à $(d)$ qui passe par $M$.
        b)~Calculer les coordonnées du projeté orthogonal $H$ de $M$ sur $(d)$, par
        substitution. c)~En déduire la distance du point $M$ à la droite $(d)$.
  \item On considère les points $A(-1\,;0)$, $B(3\,;2)$ et $C(0\,;3)$. Calculer les
        coordonnées du projeté orthogonal $K$ de $C$ sur $(AB)$, par combinaison.
\end{enumerate}

\exo[Distance d'un point à une droite]
Pour chaque cas, calculer la distance du point à la droite, en passant par son
projeté orthogonal : a)~$A(5\,;5)$ et $(d)$ : $x+y-4=0$ ; b)~$B(-1\,;9)$ et
$(d')$ : $3x-4y+14=0$ (pour le système, multiplier une équation par $3$ et l'autre
par $4$).

\exo[Raccorder une maison]
Une maison est en $M(6\,;2)$ et une route rectiligne $(d)$ a pour équation
$x-2y+8=0$ (unité : l'hectomètre). On veut raccorder la maison à la route par le
tuyau le plus court.
\begin{enumerate}
  \item Trouver une équation de la perpendiculaire à $(d)$ qui passe par $M$.
  \item Calculer les coordonnées du point $H$ de la route où arrive le tuyau.
  \item Calculer la longueur du tuyau, arrondie au mètre.
  \item Le réservoir $R(0\,;4)$ est sur la route. Calculer $MR$ et comparer.
\end{enumerate}

\partiefiche{Calcul d'un angle}

\exo[Calculer un angle]
\begin{enumerate}[label=\alph*)]
  \item Calculer la valeur exacte de l'angle des vecteurs $\vec{u}\coord{2}{0}$ et
        $\vec{v}\coord{1}{\sqrt{3}}$.
  \item On considère les points $A(1\,;1)$, $B(6\,;1)$ et $C(3\,;4)$. Calculer l'angle
        $\widehat{BAC}$, au dixième de degré.
  \item Dans le triangle $KLM$, $KL=7$, $KM=8$ et $LM=13$. Calculer l'angle $\widehat{LKM}$.
\end{enumerate}

\exo[Les trois angles d'un triangle]
On considère les points $A(-1\,;1)$, $B(3\,;-1)$ et $C(2\,;3)$. Calculer les trois angles
du triangle $ABC$, au dixième de degré, puis vérifier que leur somme vaut environ
$180^{\circ}$.

\exo[Deux expressions du produit scalaire]
On considère les points $A(1\,;2)$, $B(7\,;2)$ et $C(4\,;5)$.
\begin{enumerate}
  \item Calculer $\vect{AB}\cdot\vect{AC}$ avec les coordonnées.
  \item Calculer $AB$ et $AC$, puis écrire $\vect{AB}\cdot\vect{AC}$ avec le cosinus.
  \item En déduire la valeur exacte de l'angle $\widehat{BAC}$.
  \item Démontrer que le triangle $ABC$ est rectangle et isocèle en $C$.
\end{enumerate}

\partiefiche{Théorème d'Al-Kashi}

\exo[Calculer une longueur]
Pour chaque triangle, calculer la longueur du troisième côté :
a)~$AB=7$, $AC=15$, $\widehat{BAC}=60^{\circ}$ ; b)~$EF=5$, $EG=3$,
$\widehat{FEG}=120^{\circ}$ ; c)~$DE=\sqrt{2}$, $DF=5$, $\widehat{EDF}=45^{\circ}$ ;
d)~$RS=\sqrt{3}$, $RT=2$, $\widehat{SRT}=150^{\circ}$ ; e)~$KL=4$, $KM=6$,
$\widehat{LKM}=30^{\circ}$ (au centième).

\exo[Un angle avec Al-Kashi]
\begin{enumerate}
  \item Dans un triangle $ABC$, isoler $\cos\big(\widehat{BAC}\big)$ dans l'égalité du
        théorème d'Al-Kashi.
  \item Calculer l'angle $\widehat{BAC}$, au dixième de degré, pour $AB=5$, $AC=6$ et
        $BC=7$.
\end{enumerate}

\exo[Largeur d'un lac]
Pour mesurer la distance entre deux points $A$ et $B$ des rives d'un lac, une
géomètre se place en $S$ : elle mesure $SA=300$~m, $SB=500$~m et
$\widehat{ASB}=60^{\circ}$.
\begin{enumerate}
  \item Calculer la valeur exacte de $AB$, puis l'arrondir au mètre.
  \item Pourquoi le théorème de Pythagore ne suffit-il pas ici ?
\end{enumerate}

\exo[Diagonales et angle]
$ABCD$ est un parallélogramme avec $AB=6$, $AD=4$ et $\widehat{DAB}=60^{\circ}$.
\begin{enumerate}
  \item Calculer $\vect{AB}\cdot\vect{AD}$.
  \item Calculer la valeur exacte de la diagonale $BD$.
  \item Justifier que $\widehat{ABC}=120^{\circ}$, puis calculer la diagonale $AC$.
  \item Vérifier que $AC^{2}+BD^{2}=2\left(AB^{2}+AD^{2}\right)$.
\end{enumerate}

\partiefiche{Ensemble des points \texorpdfstring{$M$ tels que $\vect{MA}\cdot\vect{MB}=0$}{M tels que MA.MB = 0}}

\exo[Cercle de diamètre donné]
\begin{enumerate}[label=\alph*)]
  \item On considère les points $A(-3\,;1)$ et $B(5\,;-3)$. Trouver une équation de
        l'ensemble $(\mathcal{C})$ des points $M$ tels que $\vect{MA}\cdot\vect{MB}=0$.
  \item Le point $C(3\,;3)$ est-il sur $(\mathcal{C})$ ? En déduire la nature du
        triangle $ABC$.
  \item On considère les points $D(2\,;-1)$ et $E(5\,;1)$. Trouver l'ensemble
        $(\mathcal{D})$ des points $M$ tels que $\vect{DE}\cdot\vect{DM}=0$, et une
        équation de cet ensemble.
\end{enumerate}

\exo[Retrouver le cercle par le calcul]
On considère les points $A(-2\,;0)$ et $B(4\,;0)$, et un point $M(x\,;y)$.
\begin{enumerate}
  \item Calculer les coordonnées de $\vect{MA}$ et de $\vect{MB}$.
  \item Démontrer que $\vect{MA}\cdot\vect{MB}=x^{2}-2x+y^{2}-8$.
  \item En complétant le carré, démontrer que $\vect{MA}\cdot\vect{MB}=0$ équivaut à
        $(x-1)^{2}+y^{2}=9$.
  \item Reconnaître l'ensemble des points $M$ ; était-ce prévisible ?
\end{enumerate}

\exo[Angle droit sur l'axe des abscisses]
On considère les points $A(1\,;3)$ et $B(5\,;1)$, et un point $M(x\,;0)$ de l'axe des
abscisses.
\begin{enumerate}
  \item Démontrer que $\vect{MA}\cdot\vect{MB}=x^{2}-6x+8$.
  \item En déduire les points $M$ de l'axe des abscisses pour lesquels le triangle
        $AMB$ est rectangle en $M$.
  \item Trouver une équation du cercle de diamètre $[AB]$ et vérifier que ces points
        sont sur ce cercle.
  \item Trouver une équation de l'ensemble des points $N$ tels que
        $\vect{AB}\cdot\vect{AN}=0$.
\end{enumerate}

\end{multicols*}

\end{document}
